\documentclass[11pt]{article}
\usepackage{amsmath,amssymb}
\usepackage[margin=1in]{geometry}
\title{Quadrotor Dynamics, Online Reference Generation, and Neural CRaCCM Control}
\author{}
\date{}
\renewcommand{\thesubsection}{\arabic{subsection}}

\begin{document}
\maketitle

\subsection{Quadrotor dynamics and uncertainty representation}

The state is
$x=[p_x,p_y,p_z,v_x,v_y,v_z,\phi,\vartheta,\psi]^\top\in\mathbb R^9$,
where $p=[p_x,p_y,p_z]^\top$ and $v=[v_x,v_y,v_z]^\top$ are
inertial-frame positions and velocities, and $\phi,\vartheta,\psi$
are the roll, pitch, and yaw angles. The inputs
$u=[a_T,\omega_x,\omega_y,\omega_z]^\top\in\mathbb R^4$ are the
thrust acceleration and commanded body angular rates.
This is an attitude-rate model, with $a_T=T/m_{\rm nom}$ for total
thrust $T$ and nominal mass $m_{\rm nom}$.
Let $g_0=9.81~\mathrm{m\,s^{-2}}$ denote gravitational acceleration,
$R=R_z(\psi)R_y(\vartheta)R_x(\phi)$ the body-to-inertial rotation,
and $n=[n_x,n_y,n_z]^\top=Re_3$ the thrust direction.
The nominal dynamics are
\begin{align}
    f(x)&=\begin{bmatrix}
        v_x\\v_y\\v_z\\0\\0\\-g_0\\0\\0\\0
    \end{bmatrix},\qquad
    g(x)=\begin{bmatrix}
        0&0&0&0\\
        0&0&0&0\\
        0&0&0&0\\
        n_x&0&0&0\\
        n_y&0&0&0\\
        n_z&0&0&0\\
        0&1&\sin\phi\tan\vartheta&\cos\phi\tan\vartheta\\
        0&0&\cos\phi&-\sin\phi\\
        0&0&\sin\phi\sec\vartheta&\cos\phi\sec\vartheta
    \end{bmatrix}.
\end{align}
The model is used on an Euler-angle chart with $\cos\vartheta\ne0$.

For inertial-frame wind velocity
$v_{\rm w}^I(t)=[v_{{\rm w},x}(t),v_{{\rm w},y}(t),v_{{\rm w},z}(t)]^\top$,
the relative body-frame velocity is
\begin{equation}
    v_{\rm rel}(t,x)=R(x)^\top\big(v-v_{\rm w}^I(t)\big).
\end{equation}
The aerodynamic acceleration and true uncertainty are
\begin{align}
    a_{{\rm w},i}^B(t,x)&=-\beta_i v_{{\rm rel},i}(t,x)
        |v_{{\rm rel},i}(t,x)|,\qquad i\in\{x,y,z\},\\
    w(x,t)&=\begin{bmatrix}
        0_3\\
        R(x)a_{\rm w}^B(t,x)-g_0\mu n(x)+d_{\rm g}(t)\\
        w_\eta(x,t)
    \end{bmatrix},\label{eq:quad-uncertainty}
\end{align}
where $\beta_i$ are the drag coefficients divided by nominal mass,
$\mu$ is the payload coefficient in an additive approximation about
hover, $d_{\rm g}$ is an unmodeled inertial-frame acceleration disturbance,
and $w_\eta$ contains unmodeled Euler-angle-rate disturbances.
The physical plant follows
\begin{equation}\label{eq:quad-plant}
    \dot x=f(x)+g(x)u+w(x,t).
\end{equation}
In the default experiment,
$(\beta_x,\beta_y,\beta_z)=(0.09,0.23,0.07)~\mathrm{m^{-1}}$,
$v_{\rm w}^I=(0.08,-0.06,0.02)^\top~\mathrm{m\,s^{-1}}$,
$\mu=0.11$, and $\varphi_0=0.4$, with
\begin{align}
    d_{\rm g}(t)&=\begin{bmatrix}
        0.008\sin(0.71t+\varphi_0)\\
        0.008\cos(0.53t+0.3+\varphi_0)\\
        0.006\sin(0.37t+0.2+\varphi_0)
    \end{bmatrix},\\
    w_\eta(x,t)&=\begin{bmatrix}
        0.04v_{{\rm w},x}+0.003\sin(0.43t+\varphi_0)+0.002\sin\vartheta\\
        -0.04v_{{\rm w},y}+0.003\cos(0.39t+\varphi_0)+0.002\sin\phi\\
        0
    \end{bmatrix}.
\end{align}

The learned uncertainty model is
\begin{align}
    Y_\Theta(x)&=\Psi(x)\Theta,\qquad
        \Theta\in\mathbb R^{5\times4},\quad a\in\mathbb R^4,
        \label{eq:quad-representation}\\
    \Psi(x)&=\begin{bmatrix}
        0&0&0&0&0\\
        0&0&0&0&0\\
        0&0&0&0&0\\
        -v_x&0&-g_0n_x&0&0\\
        0&-v_y&-g_0n_y&0&0\\
        0&0&-g_0n_z&0&0\\
        0&0&0&1&0\\
        0&0&0&0&1\\
        0&0&0&0&0
    \end{bmatrix}.
\end{align}
At representation index $j$, the estimated uncertainty is
$Y_{\hat\Theta_j}(x)\hat a=\Psi(x)b$, where $b=\hat\Theta_j\hat a$.
Thus, the estimated model captures horizontal linear drag, a
thrust-direction acceleration bias, and roll/pitch rate biases.
The nonlinear physical disturbance need not lie exactly in this model class.

\subsection{Online generation of the desired trajectory}

The desired trajectory is generated online from prescribed position
waypoints $\{p_i^{\rm wp}\}_{i=0}^{K}\subset\mathbb R^3$ at
$T_i=i|I_k|$, with the common interval duration $|I_k|=2~\mathrm s$.
A natural cubic spline $\bar p(t)$ interpolates these positions and
supplies a waypoint-following cost guide. Let
$C_p=[I_3\;\;0_{3\times6}]$ select the position coordinates.
The complete initial reference state $x_{d,0}$ is specified separately,
with $C_px_{d,0}=p_0^{\rm wp}$, whereas intermediate attitudes and
velocities are not prescribed. Initialized at $x_d(0)=x_{d,0}$, the
desired trajectory follows the certainty-equivalent dynamics
\begin{equation}\label{eq:quad-reference}
    \dot x_d(t)=f(x_d(t))+g(x_d(t))u_d(t)
        +Y_{\hat\Theta_j}(x_d(t))\hat a(t).
\end{equation}
At each replanning time $t_\ell=\ell\Delta t_{\rm r}$, the following
certainty-equivalent MPC problem is solved:
\begin{equation}\label{eq:quad-mpc}
\begin{aligned}
    \min_{\{z_n\}_{n=0}^{N_{\rm p}},\{v_n\}_{n=0}^{N_{\rm p}-1}}\quad&
    \sum_{n=0}^{N_{\rm p}-1}\left(
        \|C_pz_n-\bar p(t_{\ell,n})\|_{Q_{\rm p}}^2
        +\|v_n-u_{\rm hov}\|_{R_{\rm p}}^2\right)\\
    &+\|C_pz_{N_{\rm p}}-\bar p(t_{\ell,N_{\rm p}})\|_{P_{\rm p}}^2\\
    \text{s.t.}\quad&z_0=x_d(t_\ell),\\
    &z_{n+1}=\Phi_{\Delta t_{\rm p}}^{\rm CE}
      (z_n,v_n;\hat\Theta_{j_\ell},\hat a(t_\ell)),\\
    &u_{\min}\le v_n\le u_{\max},\qquad n=0,\ldots,N_{\rm p}-1.
\end{aligned}
\end{equation}
Here, $t_{\ell,n}=t_\ell+n\Delta t_{\rm p}$,
$\|y\|_W^2=y^\top Wy$, $u_{\rm hov}=(g_0,0,0,0)^\top$,
$u_{\min}=(0,-2,-2,-2)^\top$, and
$u_{\max}=(20,2,2,2)^\top$. The RK4 map
$\Phi_{\Delta t_{\rm p}}^{\rm CE}$ freezes both
$\hat\Theta_{j_\ell}$ and $\hat a(t_\ell)$ during prediction.
All cost matrices are discrete-time weights, with
$Q_{\rm p}=8I_3$, $R_{\rm p}=0.03I_4$, $P_{\rm p}=60I_3$, and
$S_{\rm p}=0_{4\times4}$, disabling the optional input-increment penalty.
Waypoint following is encouraged through the running and terminal
position-guide costs, without separate waypoint penalties or intermediate
waypoint equality constraints.

The current implementation uses $N_{\rm p}=40$,
$\Delta t_{\rm p}=0.02~\mathrm s$, and
$\Delta t_{\rm r}=0.4~\mathrm s$, giving a $0.8~\mathrm s$ prediction
horizon. At each replan, the first twenty optimized reference inputs
are used successively, each held for $\Delta t_{\rm p}$.
Between replans, $x_d$ is integrated using the live $\hat a(t)$ and
current representation, without resetting it to the MPC-predicted states.
The representation-update period is $B|I_k|=8~\mathrm s$, with $B=4$;
hence $B|I_k|/\Delta t_{\rm r}=20\in\mathbb N$.
Each representation update is installed before replanning at the shared
boundary. The guide remains at the final waypoint beyond the mission end.
The default mission lasts $12~\mathrm s$ and starts from $x_{d,0}=0_9$.

\subsection{Offline training of the neural contraction metric}

The dual metric is parameterized by a neural network with features
$(v^\top,\phi,\vartheta,\psi,b^\top)^\top\in\mathbb R^{11}$,
two hidden layers of width $128$ with hyperbolic-tangent activations,
and a bias-free linear output layer of dimension $81$.
Reshaping its output into $L_\zeta(x,b)\in\mathbb R^{9\times9}$ gives
\begin{equation}\label{eq:quad-neural-metric}
    W_\zeta(x,b)=L_\zeta(x,b)^\top L_\zeta(x,b)+0.03I_9,
    \qquad M_\zeta(x,b)=W_\zeta(x,b)^{-1}.
\end{equation}
The architecture therefore guarantees $W_\zeta\succeq0.03I_9$.
The metric is independent of position, but acts on the complete
nine-dimensional differential state.

Let $F_b(x)=f(x)+\Psi(x)b$,
$A_b(x)=\partial F_b(x)/\partial x$, and let
$N(x)\in\mathbb R^{9\times5}$ be an orthonormal basis of
$\ker g(x)^\top$. For a vector $h$, write $D_xW[h]$ for the
directional state derivative with $b$ held fixed.
With contraction rate $\lambda=0.5~\mathrm{s^{-1}}$, define
\begin{align}
    C_1(x,b)&=N^\top\left(-D_xW[F_b]+A_bW+WA_b^\top
        +2\lambda W\right)N,\label{eq:quad-c1}\\
    K_i(x,b)&=D_xW[g_i]-\frac{\partial g_i}{\partial x}W
        -W\left(\frac{\partial g_i}{\partial x}\right)^\top,
        \qquad i=1,\ldots,4,\\
    C_{2,i}(x,b)&=N^\top K_i(x,b)N.\label{eq:quad-c2}
\end{align}
The training objective encourages $C_1\prec0$ and penalizes the
projected input-direction identities $C_{2,i}=0$.
For a minibatch $\mathcal B$, it is
\begin{align}
    \mathcal L(\zeta)&=1000\mathcal L_1+100\mathcal L_2
        +10\mathcal L_{\rm anchor},\\
    \mathcal L_1&=\max_{(x,b)\in\mathcal B}
        \left[\lambda_{\max}(C_1(x,b))+0.02\right]_+^2,\\
    \mathcal L_2&=\frac{1}{|\mathcal B|}
        \sum_{(x,b)\in\mathcal B}\sum_{i=1}^{4}
        \|C_{2,i}(x,b)\|_F^2,\\
    \mathcal L_{\rm anchor}&=\frac{1}{81|\mathcal B|}
        \sum_{(x,b)\in\mathcal B}
        \|W_\zeta(x,b)-W_{\zeta_0}(x,b)\|_F^2,
\end{align}
where $[r]_+=\max(r,0)$ and $\zeta_0$ denotes the supplied pretrained
neural metric. Thus, the new neural metric is obtained by fine-tuning
the supplied neural weights.

Training uses $8192$ uniformly sampled points, all $2048$ box corners
in the eleven active coordinates, and $128$
additional points near a previously identified counterexample, giving
$10368$ training points. The domain is
\begin{equation}
\begin{gathered}
    p\in[-5,5]^3,\qquad v\in[-2,2]^3,\qquad
    |\phi|,|\vartheta|\le0.35,\qquad |\psi|\le0.4,\\
    |b|\le(0.72,0.72,0.66,0.20,0.20)^\top,
\end{gathered}
\end{equation}
with componentwise inequalities. An independently sampled set of
$2048$ points is used for validation. Optimization uses Adam for
$96$ epochs with learning rate $3\times10^{-5}$, batch size $128$,
double precision, and gradient-norm clipping at $10$.
Checkpoint selection first prioritizes $C_1$ satisfaction on validation
and training boundary/counterexample points, then the margin violation,
and finally the weighted $C_2$/anchor loss; epoch $95$ is retained.
Closed-loop tracking results are not used for checkpoint selection.

After training, the neural weights are frozen. The original and retrained
neural metrics are retained as separate choices; the simulation driver
selects the original metric by default. Online updates of
$\hat\Theta_j$ and $\hat a$ do not train the neural network.
Metric derivatives are evaluated analytically, including the chain rule
\begin{equation}
    \frac{\partial W_j}{\partial\hat a_i}
    =\sum_{r=1}^{5}\frac{\partial W_\zeta}{\partial b_r}
      (\hat\Theta_j)_{ri},\qquad
    W_j(x,\hat a)=W_\zeta(x,\hat\Theta_j\hat a).
\end{equation}

\subsection{Adaptive conformal calibration and CRaCCM feedback}

The representation is initialized using $64$ observation intervals;
a separate $32$-interval dataset calibrates the learned representation
with $\hat\Theta$ frozen. Each interval contains $100$ samples at
$\Delta t=0.02~\mathrm s$. Observations are collected from the physical
plant under nominal PD excitation, and the uncertainty samples are
$w_s=\dot x_s-f(x_s)-g(x_s)u_s$.
For each completed interval, the fitted parameter and conformal score are
\begin{equation}\label{eq:quad-score}
\begin{aligned}
    a_k^{\rm fit}&\in\underset{a\in[-1,1]^4}{\arg\min}
        \sum_{t_s\in I_k}\|Y_{\hat\Theta_{j_k}}(x_s)a-w_s\|_2^2,\\
    s_k&=\max_{t_s\in I_k}
        \|Y_{\hat\Theta_{j_k}}(x_s)a_k^{\rm fit}-w_s\|_2.
\end{aligned}
\end{equation}
Every $B=4$ intervals, the representation receives a projected gradient
step on the summed fitting losses, treating the fitted parameters as fixed:
\begin{equation}
    \hat\Theta_{j+1}=\Pi_{\mathcal T}\left[
      \hat\Theta_j-10^{-4}\sum_{k\in\mathcal K_j}\sum_{t_s\in I_k}
      2\Psi(x_s)^\top
      \big(\Psi(x_s)\hat\Theta_j a_k^{\rm fit}-w_s\big)
      (a_k^{\rm fit})^\top\right].
\end{equation}
Here, $\mathcal K_j$ contains the four completed intervals in the block,
and $\Pi_{\mathcal T}$ clips each matrix entry to its prescribed trust box.
The box is centered at
\begin{equation}
    \Theta_0=\begin{bmatrix}
      0.22&0&0&0\\0.22&0&0&0\\0&0.22&0&0\\
      0&0&0.04&0\\0&0&0&0.04
    \end{bmatrix},
\end{equation}
with rowwise half-widths $(0.08,0.08,0.06,0.03,0.03)^\top$
applied to all four columns.

Let $Q_k^{\rm cp}$ be the quantile available before interval $I_k$.
The adaptive miscoverage level is updated by
\begin{equation}
    \delta_{k+1}=\delta_k+0.01
       \left(0.1-\mathbf1\{s_k>Q_k^{\rm cp}\}\right),
       \qquad\delta_0=0.1.
\end{equation}
After appending $s_k$, the next quantile is computed from the latest
$n\le32$ scores. If their order statistics are
$s_{(1)}\le\cdots\le s_{(n)}$, the implementation uses
\begin{equation}
    r=\min\!\left(n,\max\!\left(1,
       \left\lceil(1-\delta_{k+1})(n+1)\right\rceil\right)\right),
       \qquad Q_{k+1}^{\rm cp}=s_{(r)}.
\end{equation}
In particular, ranks above $n$ return the largest observed score.
These scores measure the residual of the interval-wise fitted model;
they do not directly bound the full error $w-Y_{\hat\Theta_j}\hat a$.

For the selected metric, define $M_j=W_j^{-1}$ and the geodesic energy
\begin{equation}\label{eq:quad-energy}
    E_j(x,x_d,\hat a)=\inf_{\substack{\gamma(0)=x_d\\\gamma(1)=x}}
       \int_0^1\gamma_s(s)^\top M_j(\gamma(s),\hat a)
           \gamma_s(s)\,ds.
\end{equation}
A degree-two Chebyshev path and five Clenshaw--Curtis nodes are used
in the numerical approximation.
Using the computed path and its energy, let
\begin{equation}
\begin{gathered}
    \ell_1=\gamma_s(1)^\top M_j(x,\hat a),\qquad
    \ell_0=\gamma_s(0)^\top M_j(x_d,\hat a),\qquad A=\ell_1g(x),\\
    c=\ell_1\big(f(x)+g(x)u_d+Y_{\hat\Theta_j}(x)\hat a\big)
       -\ell_0\big(f(x_d)+g(x_d)u_d
             +Y_{\hat\Theta_j}(x_d)\hat a\big)+\lambda E_j.
\end{gathered}
\end{equation}
The conformal tightening is scaled smoothly near zero control authority:
\begin{equation}
\begin{gathered}
    r_A=\min\!\left(\frac{\|A\|_2}{0.01},1\right),\qquad
    \chi(A)=r_A^3(10-15r_A+6r_A^2),\\
    \Delta_{\rm cp}=\chi(A)\|\ell_1\|_2Q_k^{\rm cp}.
\end{gathered}
\end{equation}
At each control sample, the CRaCCM correction is obtained from
\begin{equation}\label{eq:quad-craccm-qp}
\begin{aligned}
    \min_{\delta u,\,\xi\ge0}\quad&
       \frac12\delta u^\top g(x)^\top g(x)\delta u
       +\frac{10^4}{2}\xi^2\\
    \text{s.t.}\quad&A\delta u+c+\Delta_{\rm cp}\le\xi,
       \qquad u=u_d+\delta u^\star.
\end{aligned}
\end{equation}
Thus, the QP minimizes the corrective state-rate magnitude and a
contraction slack. The input bounds in~\eqref{eq:quad-mpc} constrain
$u_d$; the final corrected input $u$ has no saturation constraint in
this QP. Setting $\Delta_{\rm cp}=0$ gives the comparison without CP.

The continuous adaptation laws are
\begin{equation}\label{eq:quad-adaptation}
\begin{aligned}
    \dot{\hat a}&=\operatorname{Proj}_{\Gamma}\left(
       \hat a,\nu(\rho)\Gamma Y_{\hat\Theta_j}(x)^\top\ell_1^\top\right),\\
    \dot\rho&=-\frac{\nu(\rho)
       (\partial_{\hat a}E_j)^\top\dot{\hat a}}
       {\nu'(\rho)(E_j+\eta)},\\
    \nu(\rho)&=1.5+\frac{0.5\rho}{\sqrt{1+\rho^2}},\qquad
    \nu'(\rho)=\frac{0.5}{(1+\rho^2)^{3/2}},\qquad\eta=500.
\end{aligned}
\end{equation}
The projection is centered at zero, has outer radius $2.01$, and uses
a $0.01$ transition layer; its inner radius is $2$, which contains
$[-1,1]^4$. The initial estimates are $\hat a(0)=0_4$ and $\rho(0)=0$.
With $\Theta_{\rm c}$ denoting the pretrained representation and
$d=\Theta_{\rm c}^\top(0,0,1,0,0)^\top$, the fixed adaptation gain is
\begin{equation}
    \Gamma=3\times10^{-5}
        (\Theta_{\rm c}^\top\Theta_{\rm c}+10^{-5}I_4)^{-1}
        +9.7\times10^{-4}\frac{dd^\top}{(d^\top d)^2}.
\end{equation}
The energy derivative with respect to $\hat a$ is evaluated by quadrature
using $\partial_{\hat a_i}M_j=-M_j(\partial_{\hat a_i}W_j)M_j$.
The plant, reference, and adaptation states are integrated jointly
with RK45 over each $0.02~\mathrm s$ control hold; the computed path
is reused within the hold. In the comparisons without adaptation,
$\hat a$ and $\rho$ are held at zero and the uncertainty compensation
is omitted from the reference and feedback equations.
The interval-wise representation updates remain enabled in all four cases.

To interpret the contraction constraint, consider the ideal continuous-time
case on an interval with fixed representation and a constant parameter
$a^\star\in[-1,1]^4$ such that
$w(x,t)=Y_{\hat\Theta_j}(x)a^\star+\epsilon_j(x,t)$.
Assume an exact minimizing geodesic and the residual bound
$\|\epsilon_j(x,t)\|_2\le Q_k^{\rm cp}$ throughout that interval.
For $\tilde a=a^\star-\hat a$, define
\begin{equation}
    \mathcal V_j=\nu(\rho)(E_j+\eta)
       +\tilde a^\top\Gamma^{-1}\tilde a.
\end{equation}
The first variation of the energy, the projection inequality, and
\eqref{eq:quad-adaptation} then give
\begin{equation}\label{eq:quad-conditional-decay}
    \dot{\mathcal V}_j\le
       -2\lambda\nu(\rho)E_j+2\nu(\rho)\xi
       +2\nu(\rho)\big(1-\chi(A)\big)
            \|\ell_1\|_2Q_k^{\rm cp}.
\end{equation}
The parameter-error term cancels the modeled uncertainty mismatch,
while the $\rho$ dynamics cancel the explicit metric-parameter derivative.
When $\chi=1$ and $\xi=0$, this yields a negative tracking-energy
term under the stated assumptions.

In the implemented system, metric positivity is structural, whereas
$C_1\prec0$ and the projected $C_2$ identities are trained and checked
numerically. The $C_2$ penalty does not impose exact identities, and
the stronger full-matrix conditions $K_i=0$ are not enforced by the
neural architecture. The conformal score, finite calibration window,
online representation changes, and temporal dependence do not establish
an unconditional residual bound. Equation~\eqref{eq:quad-conditional-decay}
is therefore a conditional continuous-time interpretation, rather than
a global certificate for the sampled implementation or representation
switches. What the implemented controller directly enforces is the
relaxed, smoothly tightened QP inequality~\eqref{eq:quad-craccm-qp}
at the control samples. Across the retained twelve simulations, its
maximum recorded constraint residual is below $1.3\times10^{-14}$;
this verifies numerical QP feasibility, not zero slack or global
contraction of the uncertain physical plant.

\end{document}
